\documentclass[12pt,a4paper]{amsart}

\usepackage{amsmath,amssymb,amsthm}
\usepackage{mathtools}
\usepackage{hyperref}
\usepackage{geometry}
\geometry{margin=2.5cm}
\usepackage{booktabs}

% -------------------------------------------------------
\newtheorem{theorem}{Theorem}[section]
\newtheorem{lemma}[theorem]{Lemma}
\newtheorem{corollary}[theorem]{Corollary}
\newtheorem{proposition}[theorem]{Proposition}
\theoremstyle{definition}
\newtheorem{definition}[theorem]{Definition}
\newtheorem{remark}[theorem]{Remark}

% -------------------------------------------------------
\newcommand{\N}{\mathbb{N}}
\newcommand{\Z}{\mathbb{Z}}
\newcommand{\Q}{\mathbb{Q}}
\newcommand{\R}{\mathbb{R}}
\newcommand{\C}{\mathbb{C}}
\newcommand{\eq}[1]{e\!\left(#1\right)}

% -------------------------------------------------------
\begin{document}

\title{Waring's Problem:\\
       The Hilbert--Waring Theorem, $G(k)$ Bounds,\\
       and Wooley's Efficient Congruencing}

\author{Michael Fuhrmann}
\address{Specialist Mathematics Project, Build 80}
\email{specialist-maths@localhost}
\date{2026}

\begin{abstract}
Waring's problem (1770) asks: for each $k \geq 2$, does there exist a number
$g(k)$ such that every positive integer is a sum of at most $g(k)$ perfect
$k$-th powers?  The Hilbert--Waring theorem (1909) answers affirmatively:
$g(k)$ exists for all $k$.  The Hardy--Littlewood circle method (1920--1928)
provides an asymptotic formula for the number of representations of $n$ as
a sum of $s$ $k$-th powers, valid when $s \geq s_0(k)$.  This establishes
the existence of the ``true'' invariant $G(k)$ --- the smallest $s$ such
that every sufficiently large $n$ is a sum of $s$ $k$-th powers --- and yields
bounds on $G(k)$.

We present:  the classical values $g(2) = 4$ (Lagrange), $g(3) = 9$,
$g(4) = 19$; the distinction between $g(k)$ and $G(k)$; the Hardy--Littlewood
asymptotic formula for $r_{k,s}(n)$; the circle method derivation of
$G(k) \leq 2^k + 1$ (Weyl), $G(k) \leq k^2(k-1) + k + 3$ (Vinogradov),
and $G(k) \leq k(\log k + 4.20032)$ (Wooley, 2016, using Bourgain--Demeter--Guth
decoupling); and the sharp lower bound $G(k) \geq k + 1$.
\end{abstract}

\maketitle
\tableofcontents

% ================================================================
\section{Introduction}
% ================================================================

In 1770, Edward Waring stated in his \emph{Meditationes Algebraicae}:
\begin{quote}
  Every integer is a sum of at most 4 squares, of at most 9 cubes, of
  at most 19 fourth powers, and so on.
\end{quote}
More precisely, for each $k \geq 1$, let $g(k)$ denote the smallest $s$
such that every positive integer can be written as a sum of $s$ $k$-th
powers of non-negative integers.

\begin{definition}
  \begin{align*}
    g(k) &\;=\; \min\bigl\{s \in \N : \forall n \in \N,\; n = x_1^k + \cdots + x_s^k
                             \text{ has a solution in } x_i \geq 0\bigr\}, \\
    G(k) &\;=\; \min\bigl\{s \in \N : \forall n \text{ large enough},\;
                             n = x_1^k + \cdots + x_s^k \text{ has a solution}\bigr\}.
  \end{align*}
\end{definition}

\begin{remark}
  The distinction is important: $g(k)$ is sensitive to small values
  (e.g., the integers $1, 2^k - 1$ are ``hard''), while $G(k)$ measures
  the generic behaviour for large $n$.
  Always $G(k) \leq g(k)$.  For $k = 2$: $G(2) = g(2) = 4$.
  For $k = 3$: $G(3) \leq 7$, but $g(3) = 9$ (since $23 = 8+8+1+1+1+1+1+1+1$
  and $239$ require 9 cubes each).
\end{remark}

\textbf{Main results (known):}

\begin{center}
\begin{tabular}{ccccl}
\toprule
$k$ & $g(k)$ & $G(k)$ & $G(k)$ exact? & Source \\
\midrule
2 & 4 & 4 & Yes & Lagrange 1770 \\
3 & 9 & $\leq 7$ & No & $G(3) \geq 4$ is known \\
4 & 19 & $\leq 15$ & No & Davenport 1939 ($g(4){=}19$: Balasubramanian et al.\ 1986) \\
5 & 37 & $\leq 17$ & No & \\
6 & 73 & $\leq 24$ & No & \\
$k$ & see text & $\leq k(\log k+4.2)$ & No & Wooley 2016 \\
\bottomrule
\end{tabular}
\end{center}

% ================================================================
\section{The Case $k=2$: Lagrange's Four-Square Theorem}
\label{sec:squares}
% ================================================================

\begin{theorem}[Lagrange, 1770]\label{thm:lagrange}
  Every positive integer is a sum of four perfect squares.  That is,
  $g(2) = 4$.
\end{theorem}

The proof combines:
\begin{enumerate}
  \item \textbf{Euler's identity}: $(a^2+b^2+c^2+d^2)(e^2+f^2+g^2+h^2) = A^2+B^2+C^2+D^2$.
    This shows the set of sums of four squares is closed under multiplication.
  \item \textbf{Every prime is a sum of four squares}: For $p=2$: $2 = 1+1$.
    For odd $p$: show $a^2 + b^2 + 1 \equiv 0 \pmod p$ has a solution, use
    Minkowski's theorem (or infinite descent).
  \item \textbf{Every integer $= $ product of primes}: By (1) and (2),
    every product of primes is a sum of four squares.
\end{enumerate}

\begin{theorem}[$G(2) = 4$]\label{thm:g2}
  The value $G(2) = 4$ is tight: the integer $7 = 4+1+1+1$ requires
  exactly four squares (since $7 \equiv 7 \pmod 8$ and
  $x^2 \equiv 0, 1, 4 \pmod 8$, so $7 \not\equiv 0+0+0 \pmod 8$
  in fewer than four squares).
\end{theorem}

\begin{proof}
  Since $x^2 \in \{0, 1, 4\} \pmod 8$ for all $x$, the sum of three squares
  lies in $\{0,1,2,3,4,5,6\} \pmod 8$ (achieving all residues from $0+0+0=0$
  to $4+1+1=6$ but not $7$).  Hence any integer $\equiv 7 \pmod 8$ requires
  at least four squares.  Since $7 = 4+1+1+1$, exactly four suffice for this
  case.  Legendre's three-square theorem states that $n$ is a sum of three
  squares if and only if $n \notin \{4^a(8b+7) : a, b \geq 0\}$.  Thus
  infinitely many $n$ (e.g., $7, 15, 23, \ldots$) require four squares,
  proving $G(2) \geq 4$.  Combined with Lagrange, $G(2) = 4$.
\end{proof}

% ================================================================
\section{The Hilbert--Waring Theorem}
\label{sec:hilbert}
% ================================================================

\begin{theorem}[Hilbert--Waring, 1909]\label{thm:hilbert_waring}
  For every $k \geq 1$, $g(k)$ is finite.  That is, there exists
  $s = s(k)$ such that every positive integer is a sum of $s$ perfect
  $k$-th powers.
\end{theorem}

\begin{proof}[Proof sketch (Hilbert's method, 1909)]
  Hilbert's original proof used the theory of algebraic identities to show
  that a certain product of forms is expressible as a sum of $k$-th powers.
  The modern approach via the circle method is described in Section~\ref{sec:circle}.
  Here we sketch the Hilbert--Waring proof via Schur's lemma and Dickson's
  simpler 1939 argument.

  \textbf{Step 1.} It suffices to show that $1$ is a sum of $s$ $k$-th powers
  in $\Q$, by a density argument.  Specifically, if $1 \in k\Z_p^k$ (sums
  of $k$-th powers in $\Q_p$) for all primes $p$ and in $\R$, then a local-global
  principle (Hasse--Minkowski for quadratic forms, or an ad hoc argument) gives
  representation over $\Z$.

  \textbf{Step 2.} The key identity:
  \[
    \binom{2k}{k} \;=\; \frac{1}{k} \sum_{j=0}^{k}
    (-1)^{k-j} \binom{k}{j} (2j - k)^{2k}
    \cdot \frac{1}{k^{2k}}.
  \]
  (Hilbert's identity, proved using generating functions.  The right
  side expresses $\binom{2k}{k}$ as a rational linear combination of
  $2k$-th powers, from which the result follows by scaling.)

  The complete rigorous proof requires significant combinatorial machinery;
  see Nathanson~\cite[Chapter~7]{Nathanson1996}.
\end{proof}

% ================================================================
\section{The Circle Method for Waring's Problem}
\label{sec:circle}
% ================================================================

We follow the notation of Paper~17.  Let $M = \lfloor N^{1/k} \rfloor$ and
\[
  f(\alpha) \;=\; \sum_{m=1}^{M} \eq{m^k \alpha}.
\]
The representation count is
\[
  r_{k,s}(n) \;=\; \#\bigl\{(x_1,\ldots,x_s) \in \N_0^s :
                       x_1^k + \cdots + x_s^k = n\bigr\}
  \;=\; \int_0^1 f(\alpha)^s\, \eq{-n\alpha}\, d\alpha.
\]

\begin{remark}[Three threshold conditions for $s$]
\label{rem:thresholds}
  Three conditions on $s$ appear throughout this paper; they are
  \emph{independent and serve different roles}:
  \begin{enumerate}
    \item $s > 2k+1$: ensures absolute convergence of the singular series
      $\mathfrak{S}_{k,s}(n)$ (local condition at all primes $p$).
    \item $s > 2k$: ensures the main term $\mathfrak{S}_{k,s} \cdot \mathfrak{J}_{k,s}
      \cdot n^{s/k-1}$ dominates the minor arc contribution in the circle method.
    \item $s \leq k(k+1)/2$: the range where the Vinogradov Main Conjecture
      (Theorem~\ref{thm:bdg}) applies, giving the sharp minor arc bound
      $J_{s,k}(N) \ll N^{s+\varepsilon}$.
  \end{enumerate}
  For the Wooley bound $G(k) \leq k(\log k + 4.20032)$, all three conditions
  are satisfied simultaneously: the optimal $s \approx k \log k$ satisfies
  $s > 2k$ (for $k \geq 3$) and $s \leq k(k+1)/2$ (for $k \geq 3$), giving
  the asymptotic formula with a positive main term.
\end{remark}

\begin{theorem}[Hardy--Littlewood asymptotic formula]\label{thm:HL}
  For $s > 2k + 1$ (to ensure convergence of the singular series):
  \[
    r_{k,s}(n)
    \;=\; \mathfrak{S}_{k,s}(n)\, \mathfrak{J}_{k,s}(n)\, n^{s/k - 1}
          \;+\; O\!\left(n^{s/k - 1 - \delta}\right),
  \]
  where $\mathfrak{S}_{k,s}(n)$ is the singular series (Paper~17,
  Definition~5.2), $\mathfrak{J}_{k,s}(n) = C_{k,s} > 0$ (singular integral,
  Paper~17, Lemma~5.5), and $\delta = \delta(k,s) > 0$.
\end{theorem}

\begin{corollary}[Weyl bound for $G(k)$]\label{cor:Gk_circle}
  For all $k \geq 2$:
  \[
    G(k) \;\leq\; (k-2)\,2^{k-1} + 5.
  \]
\end{corollary}

\begin{proof}
  % BUG-B4-P19-EN-001 überprüft: Paper 17 hat \newtheorem{theorem}{Theorem}[section] und
  % \newtheorem{lemma}[theorem]{Lemma} (geteilter Zähler). Weyl's inequality ist das erste
  % Theorem/Lemma in Section 6 von Paper 17 → korrekte Nummer ist Theorem 6.1. ✓ KEIN BUG.
  Apply Weyl's inequality (Paper~17, Theorem~6.1\footnote{%
    Theorem 6.1 in Paper~17 is the Weyl inequality \texttt{thm:weyl}, the first numbered
    statement in Section~6 (``Minor Arc Bounds'') of that paper.}) with $s = (k-2)2^{k-1}+5$:
  the minor arc saving $N^{-\sigma}$ with $\sigma = 2^{1-k}$ gives a minor arc
  error of order $N^{s/k - 1 - \delta}$ for some $\delta > 0$, which is absorbed
  into the main term $N^{s/k - 1}$ of Theorem~\ref{thm:HL}.
  A careful computation of the threshold condition shows that $s = (k-2)2^{k-1}+5$
  suffices; see Vaughan~\cite[Chapter~5]{Vaughan1981} for details.
\end{proof}

% ================================================================
\section{Vinogradov's Mean Value Theorem}
\label{sec:vinogradov_mvt}
% ================================================================

The fundamental limitation of Weyl's approach is that the exponent
$2^{1-k}$ becomes tiny for large $k$.  Vinogradov (1935) proved a
dramatically better result:

\begin{theorem}[Vinogradov's Mean Value Theorem (original)]\label{thm:vmvt_original}
  For $k \geq 2$ and $s \geq 1$, let $J_{s,k}(N)$ denote the number of
  solutions to
  \[
    x_1^j + \cdots + x_s^j \;=\; y_1^j + \cdots + y_s^j \quad (j = 1, \ldots, k),
  \]
  with $1 \leq x_i, y_i \leq N$.  Equivalently,
  $J_{s,k}(N) = \int_0^1 |f(\alpha_1, \ldots, \alpha_k)|^{2s}\, d\boldsymbol\alpha$
  where $f(\boldsymbol\alpha) = \sum_{n \leq N} \eq{n\alpha_1 + \cdots + n^k \alpha_k}$.

  \textbf{Vinogradov (1935):} For $s \geq k(k+1)/2$:
  \[
    J_{s,k}(N) \;\ll\; N^{2s - k(k+1)/2 + \varepsilon}.
  \]
\end{theorem}

This bound, combined with the circle method, gives:

\begin{theorem}[Vinogradov's $G(k)$ bound]\label{thm:vinogradov_Gk}
  For $k \geq 3$:
  \[
    G(k) \;\leq\; k(\log k + \log\log k + O(1)).
  \]
\end{theorem}

The modern, sharp version was proved using the \emph{decoupling theorem}:

\begin{theorem}[Bourgain--Demeter--Guth, 2015 / Wooley, 2016]\label{thm:bdg}
  (The Vinogradov Main Conjecture.)  For $s \leq k(k+1)/2$:
  \[
    J_{s,k}(N) \;\ll\; N^{s + \varepsilon}
    \quad \text{for } s \leq k(k+1)/2.
  \]
  For $s = k(k+1)/2$ exactly:
  $J_{k(k+1)/2, k}(N) \ll N^{k(k+1)/2 + \varepsilon}$.
\end{theorem}

This was a landmark result: Wooley proved it for all $k$ using
\emph{efficient congruencing} (2016), independently of and slightly
before Bourgain--Demeter--Guth's approach via Fourier decoupling.

\begin{corollary}[Wooley's $G(k)$ bound]\label{cor:wooley}
  For all $k \geq 2$:
  \[
    G(k) \;\leq\; k(\log k + 4.20032).
  \]
\end{corollary}

\begin{proof}
  The proof uses the VMC (Theorem~\ref{thm:bdg}) in the circle method.
  We choose $s$ to minimise the effective threshold while satisfying all
  three conditions from Remark~\ref{rem:thresholds}.

  \textbf{Step 1 (Minor arc bound via VMC).}
  For any $s \leq k(k+1)/2$, the VMC gives
  $J_{s,k}(N) \ll N^{s+\varepsilon}$, which implies
  $\sup_{\alpha \in \mathfrak{m}} |f(\alpha)| \ll N^{1-1/k+\varepsilon}$
  (via Cauchy--Schwarz applied to $J_{s,k}$).  Consequently, the minor arc
  integral satisfies $|\int_{\mathfrak{m}} f^s e(-n\cdot)| \ll n^{s/k-1-\delta}$
  for some $\delta = \delta(k) > 0$ depending only on $k$.

  \textbf{Step 2 (Optimising $s$).}
  The asymptotic formula holds for any fixed $s > 2k$ (so that the
  singular series converges and the main term dominates).
  The constraint $s \leq k(k+1)/2$ is satisfied for all $s \leq k^2/2$,
  which includes the optimal choice $s = \lceil k(\log k + 4.20032)\rceil$
  since $k \log k \ll k^2/2$.

  \textbf{Step 3 (Explicit constant).}
  Wooley's 2016 analysis \cite{Wooley2016} shows that for
  $s = \lceil k(\log k + 4.20032)\rceil$, the minor arc bound is
  $O(n^{s/k-1-\delta})$ with $\delta > 0$, making the circle method effective.
  The explicit constant $4.20032$ comes from bounding the product of
  $p$-adic densities in the singular series; the method yields
  $G(k) \leq k\log k + O(k)$ with effective constants.
\end{proof}

% ================================================================
\section{Wooley's Efficient Congruencing}
\label{sec:efficient}
% ================================================================

We sketch the key idea of Wooley's efficient congruencing (2012--2016),
which proves the Vinogradov Main Conjecture and gives the best bounds on $G(k)$.

\begin{definition}[$p$-adic congruencing]
  Fix a prime $p$.  For a solution $\mathbf{x} = (x_1,\ldots,x_s)$ to
  the Vinogradov system modulo $p^l$, we say the solution has
  \emph{$p$-adic valuation} $v$ if $p^v \| x_i$ for all $i$ in some subset.
\end{definition}

\begin{theorem}[Efficient Congruencing (Wooley, 2016)]\label{thm:effcong}
  Let $k \geq 2$ and $1 \leq \tau \leq k$.  For a suitable prime $p$
  and for $s \geq \tau(k - \tau + 1) + k - \tau$:
  \[
    J_{s,k}(N) \;\leq\; J_{\tau,k}(N^{1/p})^{k/\tau}
                         \cdot N^{2s - k(k+1)/2 + \varepsilon}.
  \]
\end{theorem}

\begin{proof}[Proof idea]
  The efficient congruencing method proceeds by induction on the level $l$
  of $p$-adic approximation.  At each stage, solutions to the Vinogradov
  system modulo $p^l$ are ``lifted'' to solutions modulo $p^{l+1}$ via
  Hensel's lemma (a version adapted to the system of equations).

  The key idea is to partition the solution set according to the $p$-adic
  valuation of the variables: variables with high $p$-adic valuation
  (divisible by large powers of $p$) contribute to a \emph{nested}
  mean value theorem at a smaller scale $N/p$, while variables with
  low valuation contribute main terms bounded by combinatorial counting.

  Iterating this decomposition $k/\tau$ times and combining the resulting
  bounds via H\"older's inequality yields Theorem~\ref{thm:effcong}.  The
  full proof is over 100 pages; see~\cite{Wooley2016}.
\end{proof}

\begin{remark}[Relation to Bourgain--Demeter--Guth]
  Bourgain, Demeter, and Guth (2015) proved the same Vinogradov Main
  Conjecture using a completely different method: the \emph{$\ell^2$ decoupling
  theorem} for the moment curve.  Their proof proceeds via multilinear
  Kakeya-type estimates in harmonic analysis.  The two proofs give the
  same result but are conceptually very different: Wooley's is purely
  arithmetic (using $p$-adic structure), while BDG's is analytic
  (using oscillatory integrals and $L^p$-norms).
\end{remark}

% ================================================================
\section{Known Values and Bounds}
\label{sec:table}
% ================================================================

\subsection{The function $g(k)$}

The exact value of $g(k)$ is determined by a formula involving the
binary expansion of $2^k$:

\begin{theorem}[Formula for $g(k)$]\label{thm:gk_formula}
  Let $q = \lfloor (3/2)^k \rfloor$.  Define $s(k) = 2^k + q - 2$.
  \begin{enumerate}
    \item If $\{(3/2)^k\} + (2/3)^k < 1$ (which holds for all known $k$),
      then $g(k) = s(k)$.
    \item More precisely, $g(k) = 2^k + q - 2$ for all $k$ with
      $q \cdot 2^k + q \leq 2^k \cdot (q-1) + 2^k$, which is
      equivalent to $q(2^k - q + 2) = 2^k$.
  \end{enumerate}
\end{theorem}

This formula is proved by constructing the ``worst case'' integer:
$n = 2^k \cdot q - 1$.  This integer requires exactly $g(k)$ $k$-th powers:
it uses $(q-1)$ copies of $2^k$ and $(2^k - 1)$ copies of $1^k$.

\begin{center}
\begin{tabular}{rrrrl}
\toprule
$k$ & $g(k)$ & $q = \lfloor(3/2)^k\rfloor$ & $g(k) = 2^k + q - 2$ & Year proved \\
\midrule
2 & 4 & 2 & $4+2-2=4$ ✓ & Lagrange 1770 \\
3 & 9 & 3 & $8+3-2=9$ ✓ & Wieferich 1909 \\
4 & 19 & 5 & $16+5-2=19$ ✓ & Balasubramanian et al.\ 1986 \\
5 & 37 & 7 & $32+7-2=37$ ✓ & Chen 1958 \\
6 & 73 & 11 & $64+11-2=73$ ✓ & Pillai 1936 \\
7 & 143 & 17 & $128+17-2=143$ ✓ & Pillai 1936 \\
8 & 279 & 25 & $256+25-2=279$ ✓ & Pillai 1936 \\
\bottomrule
\end{tabular}
\end{center}

\subsection{The function $G(k)$}

\begin{center}
\begin{tabular}{rrrrl}
\toprule
$k$ & $G(k)$ exact & $G(k) \leq$ & Lower bound & Reference \\
\midrule
2 & 4 & 4 & 4 & Lagrange \\
3 & unknown & 7 & 4 & Heath-Brown 1988 \\
4 & unknown & 15 & 2 & Davenport 1939 \\
5 & unknown & 17 & 6 & Vaughan 1989 \\
6 & unknown & 24 & 9 & Vaughan 1989 \\
7 & unknown & 33 & 8 & Vaughan 1989 \\
8 & unknown & 42 & 32 & Wooley 1992 \\
large $k$ & — & $k(\log k + 4.2)$ & $k+1$ & Wooley 2016 \\
\bottomrule
\end{tabular}
\end{center}

\begin{remark}[Lower bound $G(k) \geq k+1$]
  The bound $G(k) \geq k+1$ comes from a $p$-adic obstruction:
  by Hensel's lemma and counting, the $k$-th power residues modulo $p^l$
  for appropriate $p$ show that not all integers are sums of $k$ $k$-th
  powers.  For $k = 2$: $G(2) \geq 4$ from the mod 8 argument.
  For general $k$: local obstructions at primes $p \equiv 1 \pmod k$
  give $G(k) \geq k + 1$ (see Vaughan~\cite{Vaughan1981}).
\end{remark}

% ================================================================
\section{The Singular Series for Waring's Problem}
\label{sec:singular_waring}
% ================================================================

For the $k$-th power sum $x_1^k + \cdots + x_s^k = n$, the singular
series is:
\[
  \mathfrak{S}_{k,s}(n)
  \;=\; \prod_p \sigma_p(n),
\]
where $\sigma_p(n)$ is the $p$-adic density:
\[
  \sigma_p(n) \;=\; \lim_{l \to \infty} p^{-l(s-1)}
                    \cdot \#\bigl\{\mathbf{x} \pmod{p^l} :
                           x_1^k + \cdots + x_s^k \equiv n \pmod{p^l}\bigr\}.
\]

\begin{theorem}[$\mathfrak{S}_{k,s}(n) > 0$ for large $s$]\label{thm:sing_pos}
  For $s \geq 2k + 1$, $\mathfrak{S}_{k,s}(n) > 0$ for all $n$.
  In other words, the equation $x_1^k + \cdots + x_s^k = n$ has
  $p$-adic solutions for all $p$ and all $n$.
\end{theorem}

\begin{proof}
  For the prime $p$ with $p \nmid k$: the $k$-th power map on $\Z_p$ is
  $p/(p-1)$-to-one (since $|\{x^k \pmod p\}| = (p-1)/\gcd(k,p-1) + 1$),
  so the sum of $s$ $k$-th powers covers all residues modulo $p$ for
  $s \geq 1/\text{density} + 1$.  For $p \mid k$: Hensel lifting requires
  more careful analysis, but for $s \geq 2k+1$ sufficient coverage is
  guaranteed by Waring's theorem itself applied $p$-adically.
  See Vaughan~\cite[Chapter~4]{Vaughan1981}.
\end{proof}

% ================================================================
\section{Conclusion}
% ================================================================

Waring's problem encapsulates a remarkable interplay between:
\begin{enumerate}
  \item \textbf{Global solvability} ($g(k)$: all integers, small or large),
  \item \textbf{Asymptotic solvability} ($G(k)$: all large integers),
  \item \textbf{Local solvability} (the singular series $\mathfrak{S}_{k,s}(n)$:
    $p$-adic conditions at all primes),
  \item \textbf{Analytic estimates} (exponential sums, the circle method,
    Weyl differencing, Vinogradov's mean value theorem).
\end{enumerate}

The Vinogradov Main Conjecture, now proved by Wooley (2016) and
Bourgain--Demeter--Guth (2015), gives the best possible bound
$J_{s,k}(N) \ll N^{2s-k(k+1)/2+\varepsilon}$ and consequently
$G(k) \leq k\log k + O(k)$.  Exact values of $G(k)$ for $k \geq 3$
remain unknown, representing some of the deepest open problems in
additive number theory.

% ================================================================
\begin{thebibliography}{9}

\bibitem{Waring1770}
E.~Waring.
\newblock \emph{Meditationes Algebraicae}.
\newblock Cambridge, 1770.  (3rd edition 1782.)

\bibitem{Hilbert1909}
D.~Hilbert.
\newblock Beweis für die Darstellbarkeit der ganzen Zahlen durch eine feste
  Anzahl $n$-ter Potenzen (Waringsche Problem).
\newblock \emph{Math.\ Ann.}, 67:281--300, 1909.

\bibitem{HardyLittlewood1920}
G.~H.~Hardy and J.~E.~Littlewood.
\newblock A new solution of Waring's problem.
\newblock \emph{Quart.\ J.\ Math.}, 48:272--293, 1920.

\bibitem{Vaughan1981}
R.~C.~Vaughan.
\newblock \emph{The Hardy--Littlewood Method}, 2nd~ed.
\newblock Cambridge University Press, 1997.

\bibitem{Nathanson1996}
M.~B.~Nathanson.
\newblock \emph{Additive Number Theory: The Classical Bases}.
\newblock Springer, 1996.

\bibitem{Wooley2012}
T.~D.~Wooley.
\newblock Vinogradov's mean value theorem via efficient congruencing, I.
\newblock \emph{Ann.\ of Math.\ (2)}, 175(3):1575--1627, 2012.

\bibitem{Wooley2013}
T.~D.~Wooley.
\newblock Vinogradov's mean value theorem via efficient congruencing, II.
\newblock \emph{Duke Math.\ J.}, 162(4):673--730, 2013.

\bibitem{Wooley2016}
T.~D.~Wooley.
\newblock The cubic case of the main conjecture in Vinogradov's mean value theorem.
\newblock \emph{Adv.\ Math.}, 294:532--561, 2016.

\bibitem{BDG2015}
J.~Bourgain, C.~Demeter, and L.~Guth.
\newblock Proof of the main conjecture in Vinogradov's mean value theorem for
  degrees higher than three.
\newblock \emph{Ann.\ of Math.\ (2)}, 184(2):633--682, 2016.

\bibitem{IwaniecKowalski}
H.~Iwaniec and E.~Kowalski.
\newblock \emph{Analytic Number Theory}.
\newblock American Mathematical Society, Providence, RI, 2004.

\end{thebibliography}

\end{document}
